\section*{Chapter \ref{ch:b2:banks-and-crises} --- Banks, Runs and Financial Crises}

\begin{solution}{exo:b2:banks-and-crises:1}
% figdata: decades
The 1930s (14), then the 1920s and the 2000s (13 each). None in the 1940s, 1950s and 1960s.
\end{solution}

\begin{solution}{exo:b2:banks-and-crises:2}
% figdata: paths
4.89 log points after a normal peak, 1.28 after a peak with a crisis: a gap of 3.61 points of output per person, five years on.
\end{solution}

\begin{solution}{exo:b2:banks-and-crises:3}
% code: exercises/solutions.py leverage_loss
The loss is $0.02 \times 100 = 2$: half of an equity of 4, a fifth of an equity of 10.
\end{solution}

\begin{solution}{exo:b2:banks-and-crises:4}
% ledger: b2:2008-bnp, b2:2008-northern-rock, b2:2008-aig-mmf
Investors in BNP Paribas's three funds could not withdraw: the funds stopped redemptions rather than sell at a loss, a
suspension of convertibility. Northern Rock's depositors and lenders withdrew, and the central bank lent to it. The
Reserve Primary Fund's shareholders held a claim meant to be worth one dollar a share; when it fell below, the Treasury
guaranteed money market funds three days later, an answer to withdrawals from them.
\end{solution}

\begin{solution}{exo:b2:banks-and-crises:5}
% code: exercises/solutions.py contract
With $\gamma = 1$, $R^{1/\gamma - 1} = 1$, so $c_1 = 1/(\pi + 1 - \pi) = 1$ and $c_2 = R$: the bank gives no insurance. Then
$w(f) = R(1 - f)/(1 - f) = R > 1 = c_1$ for every $f < 1$: waiting always pays more.
\end{solution}

\begin{solution}{exo:b2:banks-and-crises:6}
% code: exercises/solutions.py contract, run_threshold
$c_1 = 1/(0.5 + 0.5/\sqrt2) = 1.172$, $c_2 = \sqrt2 \times 1.172 = 1.657$, $f^* = (2 - 1.172)/1.172 = 0.707$.
\end{solution}

\begin{solution}{exo:b2:banks-and-crises:7}
% code: exercises/solutions.py contract, run_threshold
$c_1 = 1/(0.25 + 0.75 \times 1.5^{-1/2}) = 1.160$, $c_2 = 1.5^{1/2} \times 1.160 = 1.420$,
$f^* = (1.5 - 1.160)/(1.160 \times 0.5) = 0.587$. The threshold is higher than 0.561: the contract insures less, so
the bank is a little less fragile.
\end{solution}

\begin{solution}{exo:b2:banks-and-crises:8}
If the bank stops after paying $\pi c_1$, it always keeps $1 - \pi c_1$ invested, enough to pay $c_2$ to each of the
$1 - \pi$ patient depositors; waiting then pays $c_2 > c_1$ whatever others do, and no one patient runs. If the share of
impatient depositors can be larger than the bank allows for, the impatient ones who arrive after the stop are not paid when
they need to consume.
\end{solution}

\begin{solution}{exo:b2:banks-and-crises:9}
% code: exercises/solutions.py crises_and_credit
The counts by decade are those of the opening chart, 88 in all. The mean change in loans to GDP over the five years before
a crisis is 13.35 percentage points, over 63 crises with data and windows clear of the wars; it matches the second chart.
\end{solution}

\begin{solution}{exo:b2:banks-and-crises:10}
With multiplier $\lambda$ on the budget: $\pi u'(c_1) = \lambda\pi$ and $(1 - \pi)u'(c_2) = \lambda(1 - \pi)/R$, so
$u'(c_1) = Ru'(c_2)$, $c_1^{-\gamma} = Rc_2^{-\gamma}$, $c_2 = R^{1/\gamma}c_1$. Substituting in the budget,
$c_1(\pi + (1 - \pi)R^{1/\gamma - 1}) = 1$. Then $c_1 > 1$ if and only if $R^{1/\gamma - 1} < 1$, that is $1/\gamma - 1 <
0$ since $R > 1$, that is $\gamma > 1$.
\end{solution}

\begin{solution}{exo:b2:banks-and-crises:11}
$w'(f) = R\,\frac{-c_1(1 - f) + (1 - fc_1)}{(1 - f)^2} = \frac{R(1 - c_1)}{(1 - f)^2} < 0$ if $c_1 > 1$.
$w(f) \geq c_1$ is $R - Rfc_1 \geq c_1 - fc_1$, so $f \leq (R - c_1)/(c_1(R - 1))$. At $f = \pi$ the budget gives
$w(\pi) = c_2 = R^{1/\gamma}c_1 > c_1$; since $w$ is decreasing and equals $c_1$ at $f^*$, $f^* > \pi$.
\end{solution}

\begin{solution}{exo:b2:banks-and-crises:12}
Whatever share withdraws, a patient depositor who waits gets $c_2 > c_1$, so waiting is the \omterm{def:b2:static-games:best-response}{best response} and only the
impatient withdraw: $f = \pi$. The bank then has exactly $R(1 - \pi c_1) = (1 - \pi)c_2$ at date 2 and pays $c_2$ from its
own assets; the guarantee is never called upon and the tax is never collected.
\end{solution}

\begin{solution}{pb:b2:banks-and-crises:1}
% code: exercises/solutions.py BRENLAND, contract, expected_utility, run_threshold, bank_empty, wait_payoff
\textbf{Part I.}
1. $\pi = 0.25$, $R = 2$, $\gamma = 2$ (since $-1/c = c^{1-2}/(1 - 2)$).
2. Impatient 1, patient 2.
3. $0.25 \times (-1) + 0.75 \times (-0.5) = -0.625$.
4. Whether a depositor is impatient is known only to him: an insurer cannot verify a claim.
5. A patient depositor would gain by claiming to be impatient if that paid more; the contract must make the truth pay.
\textbf{Part II.}
6. $0.25c_1 + 0.75c_2/2 = 1$.
7. $c_1^{-2} = 2c_2^{-2}$, so $c_2 = \sqrt2\,c_1$.
8. $c_1 = 1/(0.25 + 0.75/\sqrt2) = 1.282$; $c_2 = 1.812$.
9. $0.25 \times 1.282 + 0.75 \times 1.812/2 = 0.320 + 0.680 = 1$.
10. $-0.25/1.282 - 0.75/1.812 = -0.609$, above $-0.625$.
11. He would get 1.282 at date 1 and could store it to date 2, against 1.812 by waiting.
\textbf{Part III.}
12. $w(f) = 2(1 - 1.282f)/(1 - f)$ while $1.282f < 1$, zero after.
13. 1.812.
14. 1.437 and 1.155.
15. $f^* = (2 - 1.282)/1.282 = 0.561$.
16. $1/1.282 = 0.780$.
17. If all others withdraw, waiting pays nothing, while withdrawing pays 1.282 to those served before the bank is empty: so
withdrawing is the \omterm{def:b2:static-games:best-response}{best response}, and the run is an equilibrium.
18. In the good equilibrium depositors receive $0.25 \times 1.282 + 0.75 \times 1.812 = 1.680$ per deposit; in the run, all
the investment is liquidated for 1. The loss is 0.680 per deposit, two-thirds of the deposits themselves.
\textbf{Part IV.}
19. $0.25 \times 1.282 = 0.320$.
20. 1.812.
21. The bank stops after paying a quarter of its depositors, so the last $1/3 - 1/4 = 1/12$ of them, impatient, get
nothing when they need it.
22. Nothing: no patient depositor runs, so the guarantee is never called upon.
23. Depositors protected whatever happens no longer care what the bank does with their money, and the bank's owners keep
the gains from risk while the guarantor bears part of the losses.
24. No: with $\gamma = 1$ the optimal contract is $c_1 = 1$, $c_2 = 2$; there is no insurance and no run.
25. The optimal first-period payout is $c_1 = 1/(\pi + (1 - \pi)R^{1/\gamma - 1})$, 1.282 here; the bank fails once a
share $1/c_1$ withdraws, 78 per cent, and waiting stops paying once $f^* = (R - c_1)/(c_1(R - 1))$ withdraw, 56 per cent.
\end{solution}
